% Self-contained note. Compile twice with pdfLaTeX.
% The bibliography is included; no external files are required.
\documentclass[11pt,a4paper]{article}
\usepackage[T1]{fontenc}
\usepackage{lmodern}
\usepackage{amsmath,amssymb,amsthm,mathtools}
\usepackage[margin=1in]{geometry}
\usepackage{microtype}
\usepackage{xcolor}
\usepackage{tikz-cd}
\tikzcdset{row sep=2.2em,column sep=2.8em}
\tikzset{
  proarrow/.style={"{\scriptscriptstyle /}" marking}
}
\usepackage[colorlinks=true,linkcolor=blue!50!black,
 citecolor=blue!50!black,urlcolor=blue!50!black]{hyperref}

\newtheorem{theorem}{Theorem}[section]
\newtheorem{proposition}[theorem]{Proposition}
\newtheorem{lemma}[theorem]{Lemma}
\newtheorem{corollary}[theorem]{Corollary}
\theoremstyle{definition}
\newtheorem{definition}[theorem]{Definition}
\theoremstyle{remark}
\newtheorem{remark}[theorem]{Remark}
\numberwithin{equation}{section}
\newcommand{\V}{\mathcal V}
\newcommand{\D}{\mathbb D}
\newcommand{\Ccal}{\mathbb C}
\newcommand{\Bcal}{\mathbb B}
\newcommand{\Kcal}{\mathbb K}
\newcommand{\Ecal}{\mathsf E}
\newcommand{\Q}{\mathsf Q}
\newcommand{\Z}{\mathbb Z}
\newcommand{\Ab}{\mathbf{Ab}}
\newcommand{\Icat}{\mathsf 1}
\newcommand{\harr}{\mathrel{\nrightarrow}}
\newcommand{\cell}{\mathrel{\Longrightarrow}}
\newcommand{\wh}{\mathbin{\odot}}
\newcommand{\emp}[1]{\varnothing_{#1}}
\newcommand{\eps}{\varepsilon}
\newcommand{\id}{\mathrm{id}}
\newcommand{\op}{\mathrm{op}}
\newcommand{\Mod}{\operatorname{Mod}}
\newcommand{\FCat}{\operatorname{FCat}}
\newcommand{\For}{\mathsf{For}}
\newcommand{\Up}{U_{\mathrm{poly}}}
\newcommand{\Ddist}{\mathbb D_{\mathrm{dist}}}
\newcommand{\FDist}{\mathsf{Dist}_{\mathrm{poly}}}

\title{Formal internal categories and distributors\\without regularity}
\author{}
\date{}

\begin{document}
\maketitle
\vspace{-1.6em}
\begin{abstract}
Cotensor products give a familiar definition of non-Cartesian internal
categories, but their existence and stability need not follow from the
ambient monoidal structure. We describe a conservative completion of
balanced bicomodule cells in which finite paths can serve as both input
and output boundaries. Category and distributor actions are then actual
cells of a specified calculus. Stable cotensors recover the usual
operations. Two examples in abelian groups explain the distinction:
finite group-like bases recover enriched categories and distributors in a
nonregular ambient category, while a small triangular category admits a
nonzero distributor action whose input path has no contextual atomic
representative. The latter action is constructed by restriction along a
formal functor, on one actual bicomodule.
\end{abstract}

\section{The question}

In Aguiar's formulation, an internal category consists of a comonoid $C$,
a $C$--$C$ bicomodule $A$, and associative, unital maps
\[
m:A\square_C A\longrightarrow A,\qquad u:C\longrightarrow A.
\]
Regularity of the monoidal category supplies the cotensors and their
compatibility with tensoring \cite{Aguiar}. Once regularity is available,
categories and distributors fit into the virtual-equipment framework.
The question here concerns the operations that remain meaningful when
that ambient condition is removed.

The usual virtual treatment of enriched distributors suggests separating
operations from represented composites \cite{CS}. There is a variance
issue, however. Before a bicomodule cotensor exists, the available data
are balanced maps \emph{into} a tensor product:
\[
N\longrightarrow M_1\otimes\cdots\otimes M_n.
\]
After orienting these as forward maps, they have one input and several
outputs. Multiplication and distributor actions have several inputs and
one output. We therefore construct a calculus with paths on both sides,
while retaining exactly the original balanced maps.

The concrete test is Theorem~\ref{thm:strong}. We construct
a formal category with carrier $A\cong\Z^5$ and a distributor with carrier
$M\cong(\Z/4)^2$. Its action $(A,M)\cell M$ is associative, unital, and
nonzero, but $(A,M)$ has no contextual representative by any actual
bicomodule. An ordinary equalizer exists; its failure to represent the
boundary in a tensor context is the point of the example.

Horizontal paths are written in traversal order. Ordinary maps and
vertical cell composites are written in functional order, so $ba=b\circ a$;
$\wh$ denotes horizontal juxtaposition. A path includes its vertices, and
$\emp C$ is the empty path at $C$. Tensor coherence maps are understood.
The general construction uses no braiding. We use $\harr$ for
bicomodules, distributors, and other horizontal proarrows, following the
slashed-arrow convention of \cite[Definition 2.1]{CS}; the slash is a
proarrow marker, not a negation. In double-category diagrams, the same
slash marks a horizontal proarrow and a double arrow denotes a cell. In diagrams whose vertices are paths,
ordinary arrows denote cells in the merge realization; commutativity
asserts equality of their composites. The category and action laws below
compare tree composites. Dashed arrows indicate unique factorizations,
except where explicitly labelled $\nexists$.

\section{A conservative calculus of paths}\label{sec:construction}

\subsection{The balanced input}

\begin{definition}[Bicomodules and equivariant maps]\label{def:bicomodule}
Write $M:C\harr D$ for a bicomodule with commuting, coassociative,
counital coactions $\lambda_M:M\to C\otimes M$ and
$\rho_M:M\to M\otimes D$. An equivariant map $F:M\to M'$ over
comonoid maps $f:C\to C'$, $g:D\to D'$ satisfies
\[
\begin{tikzcd}
M \arrow[r,"F"] \arrow[d,"\lambda_M"'] &
M' \arrow[d,"\lambda_{M'}"]\\
C\otimes M \arrow[r,"f\otimes F"'] & C'\otimes M'
\end{tikzcd}
\qquad
\begin{tikzcd}
M \arrow[r,"F"] \arrow[d,"\rho_M"'] &
M' \arrow[d,"\rho_{M'}"]\\
M\otimes D \arrow[r,"F\otimes g"'] & M'\otimes D'.
\end{tikzcd}
\]
\end{definition}

\begin{definition}[Balanced cells]\label{def:balanced}
For a path $M_i:C_{i-1}\harr C_i$, an output bicomodule
$N:D\harr E$, and comonoid maps $a:D\to C_0$, $b:E\to C_n$, a
balanced map $\phi:N\to T$, where $T=M_1\otimes\cdots\otimes M_n$,
has exterior equivariance
\begin{equation}\label{eq:exterior}
\begin{tikzcd}[column sep=2.4em]
N \arrow[r,"\phi"] \arrow[d,"\lambda_N"'] &
T \arrow[d,"\lambda_{M_1}\otimes1"]\\
D\otimes N \arrow[r,"a\otimes\phi"'] & C_0\otimes T
\end{tikzcd}
\qquad
\begin{tikzcd}[column sep=2.4em]
N \arrow[r,"\phi"] \arrow[d,"\rho_N"'] &
T \arrow[d,"1\otimes\rho_{M_n}"]\\
N\otimes E \arrow[r,"\phi\otimes b"'] & T\otimes C_n
\end{tikzcd}
\end{equation}
and interior balance. Write $T_i$ for $T$ with $C_i$ inserted between
$M_i$ and $M_{i+1}$, and put
$r_i=1\otimes\rho_{M_i}\otimes1$ and
$\ell_i=1\otimes\lambda_{M_{i+1}}\otimes1$. Balance means that the
two routes below have a common composite $\psi_i$:
\begin{equation}\label{eq:balance}
\begin{tikzcd}[column sep=3.4em]
N \arrow[r,"\phi"] \arrow[dr,"\psi_i"'] &
T \arrow[d,shift left=.7ex,"r_i"]
  \arrow[d,shift right=.7ex,"\ell_i"']\\
& T_i .
\end{tikzcd}
\qquad (1\leq i<n).
\end{equation}
For the empty path at $C$, a balanced map is $\phi:N\to I$ making
\[
\begin{tikzcd}
N \arrow[r,"\rho_N"] \arrow[d,"\lambda_N"'] &
N\otimes E \arrow[d,"\phi\otimes b"]\\
D\otimes N \arrow[r,"a\otimes\phi"'] & C
\end{tikzcd}
\]
commute. Regard these maps as cells $(M_1,\ldots,M_n)\cell N$:
\[
\begin{tikzcd}[column sep=1.8em,row sep=3.4em,
  cells={nodes={minimum width=2.4em}}]
C_0 \arrow[r,proarrow,"M_1"] \arrow[d,"a^{\op}"'] &
C_1 \arrow[r,"\cdots",""{name=top,below}] &
C_{n-1} \arrow[r,proarrow,"M_n"] &
C_n \arrow[d,"b^{\op}"]\\
D \arrow[rrr,proarrow,"N"',""{name=bottom,above}] &&& E
\arrow[Rightarrow,from=top,to=bottom,shorten <=2pt,shorten >=2pt,"\phi",
  to path={(\tikztostart) -- (\tikztostart |- \tikztotarget) \tikztonodes}]
\end{tikzcd}
\]
The displayed exterior arrows are in the opposite comonoid category.
\end{definition}

\begin{proposition}[The balanced virtual equipment]\label{prop:balanced-equipment}
For every monoidal $\V$, balanced cells form a virtual equipment
\[
\Ccal(\V)=\Mod(\V^{\op}),
\]
with comonoids as objects and \emph{opposite} comonoid maps as vertical
arrows. Its restrictions are corestrictions of coactions, and its units
are the regular bicomodules. Here $\Mod$ is the monoids-and-bimodules
construction, in its virtual-double-category form
\cite[Definition 2.8]{CS}.
\end{proposition}
\begin{proof}
Substitution is composition with the tensor of the maps inserted into
the factors. Equations~\eqref{eq:exterior}--\eqref{eq:balance}, including the empty
case, are preserved by this substitution.

Restrictions retain
the underlying object and corestrict its coactions along the actual
comonoid maps. The identity map is cartesian, since factorization through
it imposes exactly the same equivariance equations. The regular
bicomodule $C$ represents $\emp C$, with comparison given by $\eps_C$.
Indeed a balanced $\phi:N\to I$ has the unique equivariant factorization
over $(a,b)$
\[
\begin{tikzcd}
N \arrow[r,dashed,"\exists!\,\bar\phi"] \arrow[dr,"\phi"'] &
C \arrow[d,"\eps_C"]\\
& I,
\end{tikzcd}
\qquad
\bar\phi=(a\otimes\phi)\lambda_N=(\phi\otimes b)\rho_N.
\]
Coassociativity gives equivariance; applying the counit to either
equivariance square forces this formula and proves uniqueness.
In a nonempty context, insert the regular factor using the adjacent
coaction; balance identifies the two insertion formulas and the counit
removes the inserted factor. This proves the contextual unit property.
\end{proof}

\subsection{Forests and intrinsic comparisons}

\begin{definition}[The forest category]\label{def:forests}
Let $P$ be the entire globular part of a virtual equipment $\D$, retaining
arrows between different objects. Call its original horizontal arrows
\emph{atomic}. Let $\For(P)$ be the strict $2$-category with finite paths
as $1$-arrows. A $2$-cell $\Gamma\cell(K_1,\ldots,K_m)$ is a subdivision
of $\Gamma$ into $m$ consecutive, possibly empty blocks, together with an
original cell from the $j$th block to $K_j$. Cut vertices must agree.
For an empty output there is only the identity forest on the empty path.
Vertical composition substitutes forests and horizontal composition
concatenates them. In particular,
\[
\For(P)(\Gamma;K)=P(\Gamma;K)
\quad\text{for atomic }K.
\]
\end{definition}

\begin{definition}[Contextual representations]\label{def:intrinsic}
An original cell $s:\Omega\cell T$ is \emph{contextually universal} if
every original cell $\alpha$ has the unique factorization
\begin{equation}\label{eq:intrinsic}
\begin{tikzcd}[column sep=large]
(L,\Omega,R) \arrow[r,"1_L\wh s\wh1_R"] \arrow[dr,"\alpha"'] &
(L,T,R) \arrow[d,dashed,"\exists!\,\bar\alpha"]\\
& K
\end{tikzcd}
\end{equation}
for every compatible pair of paths $L,R$ and every atomic $K$.
Equivalently, substitution is a bijection
$P(L,T,R;K)\cong P(L,\Omega,R;K)$.
Let $S(P)$ be the class of these comparisons, including empty-source ones.
A path is \emph{represented} when it has such a comparison to an atom.
\end{definition}

\begin{remark}[Exterior boundaries]
The contexts in~\eqref{eq:intrinsic} may pass through other objects.
Cartesian restriction of the target gives
$\D_{f,g}(\Gamma;K)\cong P(\Gamma;K[f,g])$, so this globular test is
equivalent to the same test with arbitrary exterior vertical maps.
\end{remark}

\begin{definition}[The local merge calculus]\label{def:localization}
Define $L(P)$ by adjoining to $\For(P)$ an inverse
$s^{-1}:T\cell\Omega$ for each $s\in S(P)$, with the two inverse
equations. The equations are closed under both compositions. Equivalently,
localize at all $1_L\wh s\wh1_R$ in the hom-categories. Horizontal
composition extends because it takes such instances to such instances.
\end{definition}

\begin{theorem}[Conservativity and detection]\label{thm:conservative}
The canonical map
\begin{equation}\label{eq:conservative}
P(\Gamma;K)\longrightarrow L(P)(\Gamma;K)
\end{equation}
is a bijection for every atomic $K$, compatible with substitution.
Moreover, $\Gamma$ is isomorphic to an atom $T$ in $L(P)$ exactly when
$\Gamma$ is represented by $T$ in the sense of
Definition~\ref{def:intrinsic}.
\end{theorem}
\begin{proof}
For an atom $K$, the substitution presheaf
$\Phi_K(\Gamma)=P(\Gamma;K)$ on $\For(P)$ takes contextual comparisons
to bijections. It therefore extends to $L(P)$ by making their formal
inverses act by inverse bijections. If original cells have the same image,
apply their presheaf actions to $1_K$; this recovers the cells and proves
injectivity.

For surjectivity, express $z:\Gamma\cell K$ as a vertical sequence of
forests and inverses of contextual comparisons. Interchange provides
such a sequence. Starting with $1_K$, read it backwards. At a forest,
substitute into the current original cell. At the inverse of
$t=1_L\wh s\wh1_R$, factor the current cell uniquely as $c\,t$, using
\eqref{eq:intrinsic}; then $(c\,t)t^{-1}=c$. Each step replaces the
localized suffix by an original cell. At the end, $z$ is an original
cell. These operations are exactly the extended presheaf action, so the
answer is independent of the expression.

Finally, an isomorphism $\Gamma\to T$ is an original cell by
conservativity. Whiskering it by $L,R$ and precomposing into any atomic
output gives a bijection. Equation~\eqref{eq:conservative} identifies this
with~\eqref{eq:intrinsic}, proving detection. The converse holds by the
definition of the localization.
\end{proof}

\begin{corollary}[Universal property]\label{cor:universal}
An interpretation of $P$ in a strict $2$-category $\mathcal T$ extends
uniquely to forests. The resulting $G$ has the unique extension below
if and only if it inverts $S(P)$:
\[
\begin{tikzcd}
\For(P) \arrow[r,"\ell"] \arrow[dr,"G"'] &
L(P) \arrow[d,dashed,"\exists!\,\bar G"]\\
& \mathcal T.
\end{tikzcd}
\]
\end{corollary}
\begin{proof}
A forest must be sent to the juxtaposition of its constituent cells.
An extension across $L(P)$ must send each $s^{-1}$ to $G(s)^{-1}$;
these assignments respect exactly the defining relations of
Definition~\ref{def:localization}.
\end{proof}

\subsection{Vertical boundaries and the two-sided reduct}\label{sec:global}

\begin{lemma}[Binding comparisons]\label{lem:bindings}
For a vertical arrow $f:X\to X'$, choose
local binding presentations $f_*:X\harr X'$ and $f^*:X'\harr X$ from
restrictions of units. Their binding cells yield in $L(P)$ an adjunction
with
\[
j_f:\emp X\cell f_*f^*,\qquad
e_f:f^*f_*\cell\emp{X'}.
\]
Its triangle identities are
\[
\begin{tikzcd}[column sep=2em]
f_* \arrow[r,"j_f\wh1"] \arrow[dr,equal] &
f_*f^*f_* \arrow[d,"1\wh e_f"]\\
& f_*
\end{tikzcd}
\qquad
\begin{tikzcd}[column sep=2em]
f^* \arrow[r,"1\wh j_f"] \arrow[dr,equal] &
f^*f_*f^* \arrow[d,"e_f\wh1"]\\
& f^* .
\end{tikzcd}
\]
The base-change comparisons
\[
f_*Kg^*\cell K[f,g],\qquad
k_{f,u}:f_*u_*\cell(uf)_*,\qquad
l_{f,u}:u^*f^*\cell(uf)^*
\]
are contextually universal and hence invertible
\cite[Theorem 7.16 and Corollary 7.19]{CS}. The comparisons $k,l$ are
associative and compatible with $j,e$.
\end{lemma}
\begin{proof}
The triangle identities are the companion and conjoint binding
identities, with the represented units identified with empty paths.
The first comparison is obtained by adjoining the companion and conjoint
binding cells to $K$ and factoring through its restriction. Inserting
the other binding cells gives the inverse factorization in every
context. The same argument for pasted binding pairs gives $k,l$.
Uniqueness of binding-compatible comparisons proves their associativity
and their compatibility with $j,e$.
\end{proof}

\begin{definition}[The global merge calculus]\label{def:global}
For a square with path boundaries
\[
\begin{tikzcd}[column sep=5em,row sep=3.2em]
X \arrow[r,proarrow,"\Gamma",""{name=top,below}] \arrow[d,"f"'] &
Y \arrow[d,"g"]\\
X' \arrow[r,proarrow,"\Delta"',""{name=bottom,above}] & Y'
\arrow[Rightarrow,from=top,to=bottom,shorten <=2pt,shorten >=2pt,"a"]
\end{tikzcd}
\]
define the merge realization by
\begin{equation}\label{eq:global}
\Ecal(\D)_{f,g}(\Gamma;\Delta)
=L(P)(\Gamma;f_*\Delta g^*)
\end{equation}
in any such local presentation. More intrinsically, take the disjoint
union over the two binding presentations and identify representatives by
their unique binding-compatible isomorphisms. Only finite local choices
are used in any pasting. Normalize identity vertical maps to empty
companion and conjoint paths.

For representatives $\widehat a,\widehat b$ in~\eqref{eq:global},
horizontal pasting of adjacent squares with sides $(f,g)$ and $(g,h)$
is the outer route of
\[
\begin{tikzcd}[column sep=3em]
\Gamma\Gamma'
  \arrow[r,"\widehat a\wh\widehat b"]
  \arrow[dr,"\widehat{a\wh b}"'] &
f_*\Delta g^*g_*\Delta'h^*
  \arrow[d,"1\wh e_g\wh1"]\\
& f_*\Delta\Delta'h^* .
\end{tikzcd}
\]
If $b$ below $a$ has sides $(u,v)$ and bottom boundary $\Xi$, vertical
pasting is defined by the square
\[
\begin{tikzcd}[column sep=6em]
\Gamma \arrow[r,"\widehat a"] \arrow[d,"\widehat{ba}"'] &
f_*\Delta g^* \arrow[d,"1\wh\widehat b\wh1"]\\
(uf)_*\Xi(vg)^* &
f_*u_*\Xi v^*g^*
  \arrow[l,"k_{f,u}\wh1_\Xi\wh l_{g,v}"]
\end{tikzcd}
\]
Vertical identities are path identities; the horizontal identity on $f$
is $j_f$.
\end{definition}

\begin{proposition}[Coherence of the merge calculus]\label{prop:merge}
Definition~\ref{def:global} gives a strict double category with paths as
horizontal arrows. Its operations are independent of the local binding
presentations.
\end{proposition}
\begin{proof}
The triangle identities of Lemma~\ref{lem:bindings} prove the unit laws. Horizontal
associativity contracts disjoint counits; vertical associativity is that
of $k,l$. Interchange is the composite-counit square
\[
\begin{tikzcd}[column sep=4em]
v^*g^*g_*v_* \arrow[r,"l_{g,v}\wh k_{g,v}"]
  \arrow[d,"1\wh e_g\wh1"'] &
(vg)^*(vg)_* \arrow[d,"e_{vg}"]\\
v^*v_* \arrow[r,"e_v"'] & \emp{Y''}.
\end{tikzcd}
\]
Here $g:Y\to Y'$ and $v:Y'\to Y''$.
These diagrams commute by the same binding identities, and are
invariant under changes of local presentation by uniqueness of the
binding-compatible comparisons.
\end{proof}

\begin{theorem}[Global conservativity and boundary transport]\label{thm:global}
The inclusion of $\D$ into $\Ecal(\D)$ is bijective on cells with one
atomic output, with arbitrary exterior vertical maps. It preserves
substitution and preserves and reflects the original cartesian cells.
Companion and conjoint mates give coherent transport of both boundaries.
\end{theorem}
\begin{proof}
For one atomic output the required bijection is the top edge of
\begin{equation}\label{eq:globalconservative}
\begin{tikzcd}[column sep=3em]
\D_{f,g}(\Gamma;K)
  \arrow[r,"\cong"] \arrow[d,"\cong"'] &
\Ecal(\D)_{f,g}(\Gamma;K) \arrow[d,equal]\\
P(\Gamma;K[f,g]) \arrow[r,"\cong"'] &
L(P)(\Gamma;f_*Kg^*).
\end{tikzcd}
\end{equation}
The left bijection is cartesian factorization; the bottom one uses
Theorem~\ref{thm:conservative} and the base-change comparison. Binding
pairs cancel at internal seams, proving compatibility with substitution.
These bijections also preserve and reflect the original cartesian cells.
The mate bijections form the triangle
\[
\begin{tikzcd}[column sep=3em]
\Ecal(\D)_{f,g}(\Gamma;\Delta)
  \arrow[r,"\cong"] \arrow[dr,"\cong"'] &
L(P)(\Gamma;f_*\Delta g^*) \arrow[d,"\text{mates}","\cong"']\\
& L(P)(f^*\Gamma g_*;\Delta)
\end{tikzcd}
\]
and supply coherent transport of both boundaries. Their compatibility
with pasting follows from the binding identities of
Lemma~\ref{lem:bindings}.
\end{proof}

\begin{definition}[The two-sided balanced calculus]\label{def:Q}
The \emph{poly-virtual reduct} $\Up$ retains all cells and equalities of
this merge calculus and specifies its tree operations: connected,
acyclic rectangular pastings along atomic wires, with matching vertical
seams. Structural identity wires are contracted, unused through-wire
components are excluded, and identity wires and empty squares on
vertical maps are also retained as structural operations. Two parallel
internal wires form a cycle, so composition along an entire multiwire
boundary is a merge operation. Nullary vertices are allowed. The monad
and action equations below use only tree substitution; merge operations
remain available in the chosen realization for constructing their cells.

Finally set
\begin{equation}\label{eq:Q}
\Bcal(\V)=\Ecal(\Ccal(\V))^{\mathrm v},
\qquad
\Q(\V)=\Up\Bcal(\V).
\end{equation}
Here vertical reversal reverses both the vertical category and the
squares, retaining horizontal directions. Thus vertical arrows are now
forward comonoid maps.
\end{definition}

\begin{corollary}[Primitive and unary cells]\label{cor:unary}
The calculus $\Q(\V)$ preserves and reflects every primitive
\emph{single-input} balanced cell. Its unary squares between actual
bicomodules are exactly forward equivariant maps, and it has coherent
boundary transport.
\end{corollary}
\begin{proof}
Reverse the squares in Theorem~\ref{thm:global}. Its single-output
bijections become single-input bijections, and its two mate descriptions
exchange. The reduct retains all these cells and equalities.
\end{proof}

\section{Formal categories and cotensor recovery}\label{sec:categories}

\begin{definition}[Formal internal categories]\label{def:category}
A \emph{formal internal category} is a comonoid $C$, an actual
endobicomodule $A:C\harr C$, and globular cells
\[
\mu_A:(A,A)\cell A,\qquad \eta_A:\emp C\cell A
\]
in $\Q(\V)$, satisfying
\begin{equation}\label{eq:monadlaws}
\begin{tikzcd}[column sep=2.3em]
(A,A,A) \arrow[r,"{(\mu_A,1)}"] \arrow[d,"{(1,\mu_A)}"'] &
(A,A) \arrow[d,"\mu_A"]\\
(A,A) \arrow[r,"\mu_A"'] & A
\end{tikzcd}
\qquad
\begin{tikzcd}[column sep=2.6em]
(A,A) \arrow[dr,"\mu_A"'] &
A \arrow[l,"{(\eta_A,1)}"'] \arrow[r,"{(1,\eta_A)}"]
  \arrow[d,equal] &
(A,A) \arrow[dl,"\mu_A"]\\
& A. &
\end{tikzcd}
\end{equation}
Write $\mathsf A=(C,A,\mu_A,\eta_A)$.
\end{definition}

\begin{definition}[Formal functors]\label{def:functor}
A formal functor
$\mathsf F=(f,F):\mathsf A\to\mathsf B$ consists of a comonoid map $f$
and a unary square $F$ over $(f,f)$, satisfying
\begin{equation}\label{eq:functorlaws}
\begin{tikzcd}
(A,A) \arrow[r,"\mu_A"] \arrow[d,"{(F,F)}"'] &
A \arrow[d,"F"]\\
(B,B) \arrow[r,"\mu_B"'] & B
\end{tikzcd}
\qquad
\begin{tikzcd}
\emp C \arrow[r,"\eta_A"] \arrow[d,"{[f]}"'] &
A \arrow[d,"F"]\\
\emp D \arrow[r,"\eta_B"'] & B.
\end{tikzcd}
\end{equation}
Here $C,D$ are the object comonoids and $[f]$ is the structural empty
square on $f$.
\end{definition}

\begin{proposition}[Composition of formal functors]\label{prop:functor-category}
Formal internal categories and formal functors form a category
$\FCat(\Q(\V))$. Every formal functor has an actual equivariant
bicomodule map as its unary component.
\end{proposition}
\begin{proof}
Corollary~\ref{cor:unary} identifies the unary component with an actual
map. Unary substitution composes functors: pasting their two
diagrams~\eqref{eq:functorlaws} proves the same laws for the composite.
Identities and associativity are inherited from the cell calculus.
\end{proof}

\begin{remark}[Operations are selected cells]
These definitions select cells in the constructed collections. In
particular, multiplication is a finite expression in balanced cells and
the inverse comparisons of Section~\ref{sec:construction}; it is not an
additional generator freely adjoined to a given carrier.
\end{remark}

\begin{definition}[Stable cotensors]\label{def:stable}
For bicomodules $M:C\harr D$ and $N:D\harr E$, a cotensor cone is
an equalizer with its universal factorization of balanced maps:
\begin{equation}\label{eq:stable}
\begin{tikzcd}[column sep=3.4em,row sep=2.6em]
T \arrow[r,hook,"i"] &
M\otimes N
  \arrow[r,shift left=.7ex,"\rho_M\otimes1"]
  \arrow[r,shift right=.7ex,"1\otimes\lambda_N"'] &
M\otimes D\otimes N\\
& Z \arrow[u,"\phi"'] \arrow[ul,dashed,"\exists!\,\bar\phi"] &
\end{tikzcd}
\end{equation}
Call this cone \emph{stable} if every $K\otimes-\otimes L$ preserves it,
and write its vertex as $T=M\square_D N$.
\end{definition}

\begin{lemma}[Stable cones represent paths]\label{lem:stable}
A stable cotensor cone has unique exterior coactions making it an
equivariant bicomodule cone. Its cell $(M,N)\cell T$ in $\Ccal(\V)$
is contextually universal.
\end{lemma}
\begin{proof}
The left coaction is the unique factor in
\[
\begin{tikzcd}[column sep=3.6em]
T \arrow[r,"i"] \arrow[d,dashed,"\exists!\,\lambda_T"'] &
M\otimes N \arrow[d,"\lambda_M\otimes1"]\\
C\otimes T \arrow[r,hook,"1_C\otimes i"'] &
C\otimes M\otimes N.
\end{tikzcd}
\]
Testing coaction equations after these monic tensors of $i$ proves the
bicomodule laws. A balanced map in any context factors uniquely through
the corresponding tensor of $i$; testing its remaining equations in the
same way proves that the factor is balanced. Hence the cone is
intrinsically contextually universal in $\Ccal(\V)$.
\end{proof}

\begin{theorem}[Cotensor recovery]\label{thm:recovery}
Whenever the cotensors used in the category and functor laws are stable,
Definitions~\ref{def:category} and~\ref{def:functor} are equivalent to the
usual cotensor definitions: multiplication and unit are actual
bicomodule maps $m:A\square_C A\to A$ and $u:C\to A$, with their
cotensor associativity and unit laws, and formal functors are exactly
the corresponding changing-base functors.
\end{theorem}
\begin{proof}
By Lemma~\ref{lem:stable}, after reversal a stable cone's cell
$\iota:T\cell(M,N)$ has an inverse
$\tau:(M,N)\cell T$ in the merge realization. Precomposition with $\tau$
is a single-wire substitution along $T$, retained in $\Q(\V)$.
Consequently, whenever the relevant cotensors are stable, the category
operations correspond exactly to actual maps through
\[
\begin{tikzcd}[column sep=3em]
(A,A) \arrow[r,"\tau_{A,A}","\sim"'] \arrow[dr,"\mu"'] &
A\square_C A \arrow[d,"m"]\\
& A
\end{tikzcd}
\qquad
\begin{tikzcd}
\emp C \arrow[r,"\tau_0","\sim"'] \arrow[dr,"\eta"'] &
C \arrow[d,"u"]\\
& A,
\end{tikzcd}
\]
where $\tau_0:\emp C\cell C$ is the represented-unit comparison.
Invert the top comparisons to recover $m$ and $u$ in the merge
realization; unary conservativity identifies them with actual maps.

Iterated stable cotensors represent the same balanced cone functor.
Unique comparisons between them give the associator and unitors, and
their coherence follows because each route has the same composite to
the unparenthesized tensor. Thus~\eqref{eq:monadlaws} becomes precisely
cotensor associativity and unitality. For equivariant maps $F,G$ over
$(f,g),(g,h)$, equalizer factorization gives
\[
\begin{tikzcd}[column sep=3.5em]
M\square_D N \arrow[r,hook,"i"]
  \arrow[d,dashed,"\exists!\,F\square_g G"'] &
M\otimes N \arrow[d,"F\otimes G"]\\
M'\square_{D'}N' \arrow[r,hook,"i'"'] & M'\otimes N'.
\end{tikzcd}
\]
Cone uniqueness proves functoriality. Transporting
the diagrams~\eqref{eq:functorlaws} along the comparisons gives the
usual unit and multiplication diagrams for cotensor functors.
\end{proof}

\begin{corollary}[Regular bases and the unit base]\label{cor:regular}
For regular $\V$, formal internal categories and functors recover
Aguiar's internal categories and changing-base functors. Over the unit
comonoid, formal internal categories and functors are ordinary monoids
and monoid maps in every monoidal $\V$.
\end{corollary}
\begin{proof}
Regularity supplies the stable cotensors required by
Theorem~\ref{thm:recovery}. Over $I$, the coactions are the unit
isomorphisms, the two balance maps agree, and their equalizer is the
identity of the ordinary tensor product. This cone is always stable.
\end{proof}

\section{Distributors and their restrictions}\label{sec:distributors}

\begin{definition}[Distributors on bicomodules]\label{def:distributor}
For formal categories $\mathsf A=(C,A)$ and $\mathsf B=(D,B)$, a
distributor on an actual bicomodule $M:C\harr D$ has cells
\[
p:(A,M)\cell M,\qquad q:(M,B)\cell M
\]
satisfying the two associativity squares
\begin{equation}\label{eq:left}
\begin{tikzcd}[column sep=2.4em]
(A,A,M) \arrow[r,"{(\mu_A,1)}"] \arrow[d,"{(1,p)}"'] &
(A,M) \arrow[d,"p"]\\
(A,M) \arrow[r,"p"'] & M
\end{tikzcd}
\qquad
\begin{tikzcd}[column sep=2.4em]
(M,B,B) \arrow[r,"{(q,1)}"] \arrow[d,"{(1,\mu_B)}"'] &
(M,B) \arrow[d,"q"]\\
(M,B) \arrow[r,"q"'] & M,
\end{tikzcd}
\end{equation}
the unit triangles
\begin{equation}\label{eq:right}
\begin{tikzcd}[column sep=3em]
(A,M) \arrow[dr,"p"'] &
M \arrow[l,"{(\eta_A,1)}"'] \arrow[r,"{(1,\eta_B)}"]
  \arrow[d,equal] &
(M,B) \arrow[dl,"q"]\\
& M &
\end{tikzcd}
\end{equation}
and the commuting-action square
\begin{equation}\label{eq:commute}
\begin{tikzcd}
(A,M,B) \arrow[r,"{(p,1)}"] \arrow[d,"{(1,q)}"'] &
(M,B) \arrow[d,"q"]\\
(A,M) \arrow[r,"p"'] & M.
\end{tikzcd}
\end{equation}
These are tree equations in $\Q(\V)$. Equivariant maps over formal
functors $\mathsf F=(f,F),\mathsf G=(g,G)$ are actual bicomodule maps
$\theta:M\to M'$, with target category carriers $A',B'$, making
\[
\begin{tikzcd}
(A,M) \arrow[r,"p"] \arrow[d,"{(F,\theta)}"'] &
M \arrow[d,"\theta"]\\
(A',M') \arrow[r,"p'"'] & M'
\end{tikzcd}
\qquad
\begin{tikzcd}
(M,B) \arrow[r,"q"] \arrow[d,"{(\theta,G)}"'] &
M \arrow[d,"\theta"]\\
(M',B') \arrow[r,"q'"'] & M'.
\end{tikzcd}
\]
commute.
\end{definition}

\begin{proposition}[Recovery of cotensor distributors]\label{prop:distributor-recovery}
When the relevant category and action cotensors are stable,
Definition~\ref{def:distributor} recovers the usual cotensor bimodules
and their equivariant maps.
\end{proposition}
\begin{proof}
Use Lemma~\ref{lem:stable} to identify each represented action boundary
with its cotensor. Corollary~\ref{cor:unary} then identifies the actions
with actual maps. The same coherent comparisons used in
Theorem~\ref{thm:recovery} carry the action, unit, commutation, and
equivariance diagrams to their cotensor forms, and back.
\end{proof}

\begin{definition}[Path carriers]\label{def:carriers}
Let $\Kcal(\V)$ have one horizontal
arrow for each path of $\Bcal(\V)$ and
\[
\Kcal(\V)_{f,g}(H_1,\ldots,H_n;K)
=\Bcal(\V)_{f,g}(H_1\cdots H_n;K).
\]
Substitution is merge evaluation in $\Bcal(\V)$.
\end{definition}

\begin{proposition}[The carrier equipment]\label{prop:carriers}
The virtual double category $\Kcal(\V)$ is a virtual equipment.
Concatenation represents lists of carriers, the empty path is its unit,
and restrictions are obtained by companion and conjoint transport.
\end{proposition}
\begin{proof}
The substitution laws, concatenation comparisons, and units come from
the strict double category $\Bcal(\V)$ of Proposition~\ref{prop:merge}.
With companion and
conjoint notation now taken in the \emph{reversed} realization, boundary
bending makes
\[
K[f,g]=f_*Kg^*
\]
a cartesian restriction, with its standard square
\[
\begin{tikzcd}[column sep=5.4em,row sep=3.4em]
X \arrow[r,proarrow,"{K[f,g]}",""{name=top,below}] \arrow[d,"f"'] &
Y \arrow[d,"g"]\\
X' \arrow[r,proarrow,"K"',""{name=bottom,above}] & Y'
\arrow[Rightarrow,from=top,to=bottom,shorten <=2pt,shorten >=2pt,"\chi"]
\end{tikzcd}
\]
Its factorization property for further exterior
maps follows by combining successive companion and conjoint comparisons.
This gives all required cartesian restrictions.
\end{proof}

\begin{definition}[Formal distributors and balanced multicells]\label{def:formal-distributor}
A \emph{formal distributor} is a carrier path $H_M$ with
\eqref{eq:left}--\eqref{eq:commute}, evaluated in $\Kcal(\V)$. On a
one-bicomodule carrier this is exactly the preceding definition. A
multicell $\alpha:(M_1,\ldots,M_n)\cell N$ over $(\mathsf F,\mathsf G)$
is a cell on the carrier paths satisfying
the exterior action squares
\[
\begin{tikzcd}[column sep=2.5em]
(A,\Gamma) \arrow[r,"{(p_{M_1},1,\ldots,1)}"]
  \arrow[d,"{(F,\alpha)}"'] &
\Gamma \arrow[d,"\alpha"]\\
(A',N) \arrow[r,"p_N"'] & N
\end{tikzcd}
\qquad
\begin{tikzcd}[column sep=2.5em]
(\Gamma,B) \arrow[r,"{(1,\ldots,1,q_{M_n})}"]
  \arrow[d,"{(\alpha,G)}"'] &
\Gamma \arrow[d,"\alpha"]\\
(N,B') \arrow[r,"q_N"'] & N.
\end{tikzcd}
\]
Here $\Gamma=(M_1,\ldots,M_n)$, with carrier paths understood, and
$A,B$ and $A',B'$ are the source and target category carriers.
For an internal interface, let $\Gamma_i^+$ be $\Gamma$ with the
intermediate category carrier inserted between $M_i,M_{i+1}$.
Let $r_i$ apply $q_{M_i}$ and $\ell_i$ apply $p_{M_{i+1}}$.
The internal balance square is the first diagram below. For $n=0$,
the functors have a common domain $\mathsf A$, and the second diagram
is the required nullary balance:
\[
\begin{tikzcd}
\Gamma_i^+ \arrow[r,"r_i"] \arrow[d,"\ell_i"'] &
\Gamma \arrow[d,"\alpha"]\\
\Gamma \arrow[r,"\alpha"'] & N
\end{tikzcd}
\qquad
\begin{tikzcd}
A \arrow[r,"{(F,\alpha)}"] \arrow[d,"{(\alpha,G)}"'] &
(A',N) \arrow[d,"p_N"]\\
(N,B') \arrow[r,"q_N"'] & N.
\end{tikzcd}
\]

\end{definition}

\begin{proposition}[The distributor virtual double category]\label{prop:distributor-vdc}
Formal categories, formal functors, and the distributors and multicells
of Definition~\ref{def:formal-distributor} form the virtual double category
\[
\Ddist(\V)=
\left.\Mod(\Kcal(\V))\right|_{\FCat(\Q(\V))}.
\]
The restriction is on objects: keep the monoids whose carriers are
original atomic bicomodules, and keep all distributors between them.
Its vertical category is $\FCat(\Q(\V))$.
\end{proposition}
\begin{proof}
The monoids and monoid maps on these carriers are exactly
Definitions~\ref{def:category} and~\ref{def:functor}. The module equations
are those of Definition~\ref{def:formal-distributor}, as in
\cite[Definition 2.8]{CS}. Substitution preserves the displayed equations by moving an inserted
action through its block or across an adjacent interface. Empty blocks
use the nullary equation. Identities and substitution laws are
inherited from $\Kcal(\V)$.
\end{proof}

\begin{theorem}[Distributor restrictions and units]\label{thm:restrictions}
The virtual double category $\Ddist(\V)$ is a virtual equipment.
Every distributor has restrictions along all formal functors. The
regular distributor $(A,\mu_A,\mu_A)$ is the unit at $\mathsf A$, with
universal cell $\eta_A$.
\end{theorem}
\begin{proof}
Given $N$ and formal
functors $\mathsf F=(f,F),\mathsf G=(g,G)$ into its endpoints, let
$\chi:H=f_*H_Ng^*\cell H_N$ be the ambient cartesian cell. Define actions
by its unique factorizations below, writing $N$ for its carrier $H_N$:
\begin{equation}\label{eq:restrictactions}
\begin{tikzcd}
(A,H) \arrow[r,dashed,"\exists!\,p_H"]
  \arrow[d,"{(F,\chi)}"'] &
H \arrow[d,"\chi"]\\
(A',N) \arrow[r,"p_N"'] & N
\end{tikzcd}
\qquad
\begin{tikzcd}
(H,B) \arrow[r,dashed,"\exists!\,q_H"]
  \arrow[d,"{(\chi,G)}"'] &
H \arrow[d,"\chi"]\\
(N,B') \arrow[r,"q_N"'] & N.
\end{tikzcd}
\end{equation}
The following pasting tests associativity after postcomposition with $\chi$:
\[
\begin{tikzcd}[column sep=2.5em,row sep=2.6em]
(A,A,H) \arrow[r,"{(\mu_A,1)}"] \arrow[d,"{(F,F,\chi)}"'] &
(A,H) \arrow[r,"p_H"] \arrow[d,"{(F,\chi)}"] &
H \arrow[d,"\chi"]\\
(A',A',N) \arrow[r,"{(\mu_{A'},1)}"]
  \arrow[d,"{(1,p_N)}"'] &
(A',N) \arrow[r,"p_N"] \arrow[d,"p_N"] &
N \arrow[d,equal]\\
(A',N) \arrow[r,"p_N"'] & N \arrow[r,equal] & N.
\end{tikzcd}
\]
Using~\eqref{eq:restrictactions} on the left route identifies the two
outer composites with the images under $\chi$ of the desired action
composites. Cartesian uniqueness makes them equal.
The unit, right action, and commuting equations follow by the same
factorization argument. A factor of a balanced multicell through $\chi$
is balanced because postcomposition with $\chi$ reflects its defining
equations. Thus $\chi$ is cartesian also for distributor multicells.

The regular distributor $(A,\mu_A,\mu_A)$ is the unit at $\mathsf A$,
with universal cell $\eta_A$. Its contextual factorization inserts an
adjacent action; balance identifies the insertion formulas and the
action unit laws give their inverse. In the empty context the formula is
$p_N(F,\alpha)=q_N(\alpha,G)$. Consequently $\Ddist(\V)$ is a virtual
equipment, with restrictions and units constructed above.
\end{proof}

\begin{corollary}[The two-sided distributor calculus]\label{cor:distributor-completion}
Define
\[
\FDist(\V)=\Up\Ecal(\Ddist(\V)).
\]
There is no reversal at this stage. Its atomic arrows are individual
distributors, each of which has its own bicomodule-path carrier. It
preserves and reflects all balanced single-output distributor multicells,
and formal functors have companions and conjoints obtained by restricting
regular distributors.
\end{corollary}
\begin{proof}
Apply Theorem~\ref{thm:global} to the virtual equipment supplied by
Theorem~\ref{thm:restrictions}, and take the poly-virtual reduct without
reversal. Lemma~\ref{lem:bindings} supplies the companions and conjoints.
This second localization uses distributor balancing; its representing
comparisons belong to that new calculus.
\end{proof}

\section{Recovery in a nonregular base}\label{sec:old}

\begin{proposition}[Ambient nonregularity]\label{prop:ab-nonregular}
The category $(\Ab,\otimes_{\Z},\Z)$ is not regular in Aguiar's sense:
tensoring with $\Z/2$ can fail to preserve a coreflexive equalizer.
\end{proposition}
\begin{proof}
The coreflexive pair
\[
r,s:\Z\rightrightarrows\Z\oplus\Z,\qquad
r(n)=(n,2n),\quad s(n)=(n,0)
\]
has equalizer zero. Put $R=\Z/2$. After tensoring, $\bar r=\bar s$, so
the comparison from the tensored cone to the new equalizer is
\[
\begin{tikzcd}[column sep=3em]
0 \arrow[r] \arrow[d,"\kappa"'] &
R \arrow[r,shift left=.7ex,"\bar r"]
  \arrow[r,shift right=.7ex,"\bar s"'] \arrow[d,equal] &
R\oplus R \arrow[d,equal]\\
R \arrow[r,equal] &
R \arrow[r,shift left=.7ex,"\bar r"]
  \arrow[r,shift right=.7ex,"\bar s"'] &
R\oplus R.
\end{tikzcd}
\]
The bottom row is an equalizer and $\kappa$ is not an isomorphism.
Projection onto the first coordinate is a common retraction of $r,s$;
even this coreflexive equalizer is not preserved.
\end{proof}

\begin{proposition}[Stable matrix cotensors]\label{prop:matrix}
For a finite set $S$, let
\[
C_S=\bigoplus_{s\in S}\Z e_s,\qquad
\delta(e_s)=e_s\otimes e_s,\quad\eps(e_s)=1
\]
be its group-like comonoid. Bicomodules $C_S\harr C_T$ are matrices
$(M_{st})$ of abelian groups. Their cotensors are stable and given by
\begin{equation}\label{eq:matrix}
(M\square_{C_T}N)_{su}
=\bigoplus_{t\in T}M_{st}\otimes N_{tu}.
\end{equation}
\end{proposition}
\begin{proof}
The coordinate endomorphisms of a coaction are orthogonal idempotents
summing to the identity. Commutation of the two coactions gives the
matrix decomposition; conversely, the matrix entries define the two
coactions. The two balance maps insert the two intermediate basis elements.
On matching indices they agree; on mismatching indices a coordinate
projection of their difference is the identity. This calculation survives
arbitrary tensoring, so these equalizer cones are stable, with no
flatness requirement on the entries.
\end{proof}

\begin{corollary}[Enriched categories]\label{cor:enriched-categories}
Formal categories over finite group-like comonoids, and formal functors
between them, recover finite-object $\Ab$-enriched categories and
enriched functors.
\end{corollary}
\begin{proof}
Proposition~\ref{prop:matrix} and Theorem~\ref{thm:recovery} identify the
component maps as $A_{st}\otimes A_{tu}\to A_{su}$ and
$\Z\to A_{ss}$. Comonoid maps between the $C_S$'s are functions of
indexing sets: a group-like element over $\Z$ has exactly one coefficient
$1$ and all others $0$. Thus formal functors recover enriched functors.
\end{proof}

\begin{proposition}[Enriched distributors and composition]\label{prop:enriched-distributors}
Apply the distributor construction to the sector whose carrier paths
have only finite group-like bases. Its formal distributors are, up to
their representing comparisons, the usual enriched distributors.
Distributor composition is represented by the relative tensor product
\[
(M\otimes_{\mathsf B}N)_{su}
=
\frac{\displaystyle\bigoplus_t M_{st}\otimes N_{tu}}
{\bigl\langle (mb)\otimes n-m\otimes(bn)\bigr\rangle}.
\]
Here $b\in B_{tt'}$, $m\in M_{st}$ and $n\in N_{t'u}$, so the two
terms lie in the $t'$ and $t$ summands.
\end{proposition}
\begin{proof}
Every carrier path has the stable matrix representative of
Proposition~\ref{prop:matrix}. Transporting its actions along that
comparison gives exactly an enriched distributor, by
Proposition~\ref{prop:distributor-recovery}.
Tensor products and finite direct sums preserve the displayed action
coequalizer. More explicitly, write
$X=\bigoplus_{t,t'}M_{st}\otimes B_{tt'}\otimes N_{t'u}$,
$Y=\bigoplus_t M_{st}\otimes N_{tu}$ and
$Q=(M\otimes_{\mathsf B}N)_{su}$. Balanced maps have the factorization
\[
\begin{tikzcd}[column sep=3em]
X \arrow[r,shift left=.7ex,"q_M\otimes1"]
  \arrow[r,shift right=.7ex,"1\otimes p_N"'] &
Y \arrow[r,two heads,"\pi"] \arrow[dr,"\alpha"'] &
Q \arrow[d,dashed,"\exists!\,\bar\alpha"]\\
&& Z.
\end{tikzcd}
\]
The same diagram works in every tensor context. Hence the quotient has
the contextual factorization property for balanced distributor cells.
Iterated representatives are coherent by uniqueness.
\end{proof}

\begin{remark}[What this example establishes]
Proposition~\ref{prop:ab-nonregular} concerns the ambient category;
Proposition~\ref{prop:matrix} shows that all cotensors in this sector
are stable. The next example instead places the failure of stability
inside a distributor action boundary.
\end{remark}

\section{A distributor with an unrepresented action boundary}
\label{sec:strong}

The construction has two steps. First form a category whose own
composition paths are stable. Then restrict a stable distributor action
along a formal functor. The resulting action remains a well-defined
cell, while its input path loses contextual representability.

\subsection{The category and the functor}

\begin{definition}[Triangular data]\label{def:triangular-data}
Let $D=\Z g\oplus\Z x$ with
\[
\delta(g)=g\otimes g,\qquad
\delta(x)=g\otimes x+x\otimes g,\qquad
\eps(g)=1,\quad\eps(x)=0.
\]
Set $E=\Z e\oplus D$, where $e$ is group-like. Let
$H=\Z a\oplus\Z b$ be the $E$--$E$ bicomodule
\begin{equation}\label{eq:H}
\lambda_H(h)=e\otimes h,\qquad
\rho_H(a)=a\otimes g,\qquad
\rho_H(b)=b\otimes g+2a\otimes x.
\end{equation}
Also let $J=\Z t$ with $\lambda_J(t)=e\otimes t$ and
$\rho_J(t)=t\otimes g$, and define $\varphi:H\to J$ by
$\varphi(a)=0$, $\varphi(b)=t$.
\end{definition}

\begin{lemma}[Validity of the triangular data]\label{lem:triangular-data}
The displayed formulas make $E$ a comonoid, $H,J$ bicomodules over
$E$, and $\varphi$ an equivariant bicomodule map.
\end{lemma}
\begin{proof}
The comonoid laws for $E$ hold on $e,g,x$: the two iterated coproducts
of $x$ are the sum of the three tensors with one $x$ and two $g$'s,
and the counit laws follow from $\eps(g)=1$, $\eps(x)=0$.
For the right coaction on $H$, write
$\rho_H(h)=h\otimes g+d(h)\otimes x$ with $d(a)=0$, $d(b)=2a$.
The equation $d^2=0$ proves coassociativity, the first square below.
Counitality and commuting coactions follow from the formulas; the
laws for $J$ use that $e,g$ are group-like.
Right equivariance of $\varphi$ is the second square, since
$\varphi d=0$; left equivariance follows from the common support $e$:
\[
\begin{tikzcd}[column sep=2.1em]
H \arrow[r,"\rho_H"] \arrow[d,"\rho_H"'] &
H\otimes E \arrow[d,"1\otimes\delta_E"]\\
H\otimes E \arrow[r,"\rho_H\otimes1"'] & H\otimes E\otimes E
\end{tikzcd}
\qquad
\begin{tikzcd}[column sep=2.1em]
H \arrow[r,"\varphi"] \arrow[d,"\rho_H"'] &
J \arrow[d,"\rho_J"]\\
H\otimes E \arrow[r,"\varphi\otimes1"'] & J\otimes E.
\end{tikzcd}
\]

\end{proof}

\begin{lemma}[Triangular categories]\label{lem:triangular}
For $K=H$ or $J$, the carrier $A_K=E\oplus K$ defines a formal category,
with unit the inclusion of $E$ and multiplication the fold
\[
m_K:E\oplus K\oplus K\longrightarrow E\oplus K,\qquad
(c,k_1,k_2)\longmapsto(c,k_1+k_2).
\]
Every finite composition path is stably represented; for $n\geq1$,
\[
A_K^{\square_E n}\cong E\oplus K^{\oplus n}.
\]
\end{lemma}
\begin{proof}
Expand $(E\oplus K)\otimes(E\oplus K)$ into its four blocks.
The $E,E$, $E,K$, and $K,E$ blocks have the stable unit cotensors
$E$, $K$, and $K$. The $K,K$ block has stable equalizer zero:
projection of the inserted $E$ factor onto its $e$ coordinate takes
the two balance maps to zero and the identity, respectively. Thus every
equalizing map is zero, also after tensoring.

For a longer word, remove regular $E$ factors using their stable unit
cones. A block with no $K$ contributes $E$, one $K$ contributes $K$,
and two or more $K$'s contribute zero. This proves the asserted
representatives and stability. Associativity is the fold square
\[
\begin{tikzcd}[column sep=4em]
E\oplus K^{\oplus3} \arrow[r,"\nabla_{12}"]
  \arrow[d,"\nabla_{23}"'] &
E\oplus K^{\oplus2} \arrow[d,"m_K"]\\
E\oplus K^{\oplus2} \arrow[r,"m_K"'] & E\oplus K.
\end{tikzcd}
\]
Here $\nabla_{12}$ and $\nabla_{23}$ fold the indicated pair of $K$
summands; both routes add all three. The unit
equations are the summand inclusions. Theorem~\ref{thm:recovery}
converts these into formal category laws.
\end{proof}

\begin{corollary}[The restriction functor]\label{cor:triangular-functor}
Write $A=E\oplus H$ and $B=E\oplus J$, with the resulting categories
$\mathsf A,\mathsf B$. The map
\[
F=1_E\oplus\varphi:A\longrightarrow B
\]
defines a formal functor $\mathsf F=(1_E,F):\mathsf A\to\mathsf B$.
\end{corollary}
\begin{proof}
By Lemma~\ref{lem:triangular-data}, $F$ is equivariant. It preserves
the unit and the folds of Lemma~\ref{lem:triangular} through
\[
\begin{tikzcd}[column sep=3.8em]
E\oplus H^{\oplus2} \arrow[r,"m_H"]
  \arrow[d,"1_E\oplus\varphi^{\oplus2}"'] &
A \arrow[d,"F"]\\
E\oplus J^{\oplus2} \arrow[r,"m_J"'] & B
\end{tikzcd}
\qquad
\begin{tikzcd}
E \arrow[r,hook,"u_A"] \arrow[dr,"u_B"'] &
A \arrow[d,"F"]\\
& B.
\end{tikzcd}
\]
Theorem~\ref{thm:recovery} gives the formal functor laws.
The carriers $A,B,E$ have ranks five, four, and three, respectively.
\end{proof}

\subsection{An actual bicomodule with a nonzero formal action}

\begin{definition}[The distributor carrier]\label{def:M}
Put $N=\Z/4$ and $M=Nu\oplus Nv$, with
\[
\lambda_M(u)=e\otimes u,\qquad
\lambda_M(v)=g\otimes v,
\]
and the right coaction of the unit comonoid $\Z$.
Let $\Icat$ denote the identity formal category at $\Z$.
\end{definition}

\begin{proposition}[The stable action]\label{prop:stable-action}
The carrier $M$ admits a $\mathsf B$--$\Icat$ distributor structure.
Its action paths $(B,M)$ and $(B,B,M)$ are stably represented by
$M\oplus N$ and $M\oplus N\oplus N$, respectively. Its left action
is induced by the actual equivariant map
\[
p_{B,0}:M\oplus N\longrightarrow M,\qquad
((ru+sv),w)\longmapsto(r+w)u+sv.
\]
\end{proposition}
\begin{proof}
For $(B,M)$, the regular $E$ block gives $M$;
$J\otimes Nu$ has mismatched supports and gives zero; and both
coactions on $J\otimes Nv$ insert $g$, giving all of $N$.
These computations persist under arbitrary tensoring.
The extra $N$ summand and its image $Nu$ both have left support $e$,
so this is equivariant. The triple $(B,B,M)$ is stably represented by
$M\oplus N\oplus N$, and the action square becomes
\[
\begin{tikzcd}[column sep=4em]
M\oplus N\oplus N \arrow[r,"1_M\oplus+"] \arrow[d,"\sigma"'] &
M\oplus N \arrow[d,"p_{B,0}"]\\
M\oplus N \arrow[r,"p_{B,0}"'] & M.
\end{tikzcd}
\]
Here $\sigma(m,w_1,w_2)=(p_{B,0}(m,w_2),w_1)$; both routes send
$((ru+sv),w_1,w_2)$ to $(r+w_1+w_2)u+sv$. The unit selects $M$.
Hence this gives an associative, unital formal action $p_B$.
The right action $q:(M,\Z)\cell M$ is the unit comparison and commutes
with $p_B$ by its coherence.
\end{proof}

\begin{proposition}[Restriction on the actual carrier]\label{prop:restricted-action}
The same bicomodule $M$ admits a $\mathsf A$--$\Icat$ distributor
structure, obtained by restriction along
$\mathsf F$ of Corollary~\ref{cor:triangular-functor}. Its left action
is defined by the outer route
\begin{equation}\label{eq:pA}
\begin{tikzcd}[column sep=4em]
(A,M) \arrow[r,"{(F,1_M)}"] \arrow[d,"p_A"'] &
(B,M) \arrow[d,"\tau_{B,M}","\sim"']\\
M & M\oplus N \arrow[l,"p_{B,0}"]
\end{tikzcd}
\end{equation}
where $\tau_{B,M}$ is the stable comparison already constructed.
Equivalently, $p_A=p_B(F,1_M)$, a tree substitution in $\Q(\Ab)$.
Its right action is the unit action $q$ of
Proposition~\ref{prop:stable-action}.
\end{proposition}
\begin{proof}
Since the base map of $\mathsf F$ is $1_E$, restriction retains the
actual carrier $M$. Associativity follows by pasting
\[
\begin{tikzcd}[column sep=2.5em,row sep=2.6em]
(A,A,M) \arrow[r,"{(\mu_A,1)}"] \arrow[d,"{(F,F,1)}"'] &
(A,M) \arrow[r,"p_A"] \arrow[d,"{(F,1)}"] &
M \arrow[d,equal]\\
(B,B,M) \arrow[r,"{(\mu_B,1)}"] \arrow[d,"{(1,p_B)}"'] &
(B,M) \arrow[r,"p_B"] \arrow[d,"p_B"] &
M \arrow[d,equal]\\
(B,M) \arrow[r,"p_B"'] & M \arrow[r,equal] & M.
\end{tikzcd}
\]
The two routes around the outside are the two $A$-action composites,
after using the definition of $p_A$ on the left route.
Unitality is the left triangle in
\[
\begin{tikzcd}[column sep=3em]
M \arrow[r,"{(\eta_A,1)}"] \arrow[dr,equal] &
(A,M) \arrow[r,"{(F,1)}"] \arrow[d,"p_A" description] &
(B,M) \arrow[dl,"p_B"]\\
& M &
\end{tikzcd}
\]
whose outer triangle is the functor unit law followed by the $B$-action
unit law.
The commuting-action equation follows by substituting $F$ into that for
$p_B,q$. Thus $(M,p_A,q)$ is a formal distributor
$\mathsf A\harr\Icat$ on one actual bicomodule. The construction of
$p_A$ uses the inverse comparison for the stable path $(B,M)$.
\end{proof}

\subsection{The unstable cone and the contextual obstruction}

\begin{lemma}[The unstable action cotensor]\label{lem:unstable-cotensor}
The ordinary equalizer $T=A\square_E M$ exists and carries an
equivariant $E$--$\Z$ bicomodule cone. After tensoring with $R=\Z/2$,
the comparison from this cone to the new equalizer is neither monic
nor epic. In particular, the action cotensor is not stable.
\end{lemma}
\begin{proof}
The regular block of $T$ is
$M$. The $H\otimes Nu$ block is zero by its supports, while an element
\[
a\otimes\alpha v+b\otimes\beta v\in H\otimes Nv
\]
is balanced exactly when $2\beta=0$ in $N$. Indeed its balance
obstruction is $2a\otimes x\otimes\beta v$. Consequently
\begin{equation}\label{eq:badcotensor}
T\cong M\oplus Na\oplus N[2]b
\cong(\Z/4)^3\oplus\Z/2.
\end{equation}
Writing $z$ for the last, order-two generator, its cone
$i:T\hookrightarrow A\otimes M$ is
\begin{equation}\label{eq:badcone}
u\longmapsto e\otimes u,\quad
v\longmapsto g\otimes v,\quad
a\longmapsto a\otimes v,\quad
z\longmapsto2b\otimes v.
\end{equation}
Give $u,a,z$ left support $e$, and $v$ left support $g$; this makes
$T$ an $E$--$\Z$ bicomodule and $i$ an equivariant cone.

After tensoring with $R=\Z/2$, the nonzero class
$z\otimes\overline1$ maps to zero. Meanwhile
\begin{equation}\label{eq:newbalanced}
b\otimes v\otimes\overline1
\end{equation}
becomes balanced: its obstruction is a multiple of $2$. This element
cannot lie in the image of $i\otimes1_R$, whose $b$ coordinate is zero.
Thus the tensored cone fails both injectivity and surjectivity onto
the new equalizer $T_R^{\mathrm{bal}}$. The comparison is
\[
\begin{tikzcd}[column sep=4em]
T\otimes R \arrow[r,"i\otimes1_R"] \arrow[d,"\kappa"'] &
A\otimes M\otimes R \arrow[d,equal]\\
T_R^{\mathrm{bal}} \arrow[r,hook,"i_R"'] & A\otimes M\otimes R.
\end{tikzcd}
\]
Both objects on the left have underlying group $(\Z/2)^4$.
In the source basis $(u,v,a,z)$ and the target basis given by
$e\otimes u,g\otimes v,a\otimes v,b\otimes v$ after reduction,
$\kappa$ has matrix $\operatorname{diag}(1,1,1,0)$. The failure is in
this canonical cone, whose image has dimension three.
\end{proof}

\begin{theorem}[An unrepresented action boundary]\label{thm:strong}\label{prop:strong}
The action boundary $(A,M)$ has no atomic contextual representative.
Its formal action operates nontrivially on the new balanced element
\eqref{eq:newbalanced}.
\end{theorem}
\begin{proof}
In the original globular calculus $P$ of $\Ccal(\Ab)$, the cone $i$
is a cell $(A,M)\cell T$. It represents the empty-context balanced-map
functor: factor a balanced map through its ordinary equalizer, and
test equivariance of the factor after $1_E\otimes i$. This map is
monic because $E$ is free. The right coaction is the unit coaction.

Let $W=\Z/2$ have left $E$-coaction $w\mapsto e\otimes w$ and right
unit coaction. With $R$ regarded as a $\Z$--$\Z$ bicomodule, the map
\[
\beta:W\longrightarrow A\otimes M\otimes R,\qquad
\overline1\longmapsto b\otimes v\otimes\overline1
\]
satisfies both interior balances and exterior equivariance. It is a
cell in $P(A,M,R;W)$ for which the factorization triangle fails:
\[
\begin{tikzcd}[column sep=4.4em,row sep=2.6em]
W \arrow[r,"\beta"] \arrow[d,dashed,"\nexists\,h"'] &
A\otimes M\otimes R\\
T\otimes R \arrow[ur,"i\otimes1_R"'] &
\end{tikzcd}
\]
The dashed arrow denotes the nonexistent factor: the $b$ coordinate
of $\beta(\overline1)$ is nonzero, while that of every element in the
image of $i\otimes1_R$ vanishes.
Hence $i$ fails the contextual test of Definition~\ref{def:intrinsic}.

Every other contextual representative would first represent the same
empty-context functor, so its cone would be $i$ composed with a
bicomodule isomorphism:
\[
\begin{tikzcd}
S \arrow[r,"\theta","\sim"'] \arrow[dr,"j"'] &
T \arrow[d,"i"]\\
& A\otimes M.
\end{tikzcd}
\]
Here $j$ is the proposed representing cone. Tensoring $\theta$ with $R$ cannot
change this failure of factorization. No alternative atom can therefore
represent the path. By Theorem~\ref{thm:conservative}, $(A,M)$ is not
isomorphic to any atom in the globular localization, or in its reversed
merge realization.

In the forward realization, $\beta$ is the primitive cell
$W\cell(A,M,R)$. Applying $F$ to its first output replaces $b$ by $t$.
The stable comparison for $(B,M)$ followed by $p_B$ then sends
$t\otimes v\otimes\overline1$ to $u\otimes\overline1$. Thus, in the
merge realization, the entire evaluation is
\[
\begin{tikzcd}[column sep=2.8em,row sep=3em]
W \arrow[r,"\beta"] \arrow[d,"\gamma"'] &
(A,M,R) \arrow[r,"{(F,1,1)}"] \arrow[d,"p_A\wh1_R"] &
(B,M,R) \arrow[d,"\tau_{B,M}\wh1_R","\sim"']\\
(M,R) \arrow[r,equal] &
(M,R) &
(M\oplus N,R) \arrow[l,"p_{B,0}\wh1_R"]
\end{tikzcd}
\]
Here $(M,R)$ is stably represented by $M\otimes R$.
Under this comparison, unary conservativity identifies $\gamma$ with
the actual map $\overline1\mapsto u\otimes\overline1\neq0$.
The reduct retains the resulting cells and equality.
\end{proof}

\begin{corollary}[The action on the ordinary cotensor]\label{cor:ordinary-action}
The map $p_T:T\to M$ induced by the formal action is nonzero on the
order-two summand of $T$: it sends its generator $z$ to $2u$.
\end{corollary}
\begin{proof}
The primitive cone gives a triangle in the merge realization
\[
\begin{tikzcd}
T \arrow[r,"\iota"] \arrow[dr,"p_T"'] &
(A,M) \arrow[d,"p_A"]\\
& M,
\end{tikzcd}
\]
where unary conservativity identifies $p_T$ with the actual map
\[
p_T:T\longrightarrow M,\qquad
(ru+sv,\alpha,\beta)\longmapsto(r+\beta)u+sv,
\qquad \beta\in N[2],
\]
so $p_T(z)=2u\neq0$. Nevertheless, the contextual action above cannot
be evaluated by first factoring through $T\otimes R$: that
factorization does not exist by Theorem~\ref{thm:strong}.
\end{proof}

\begin{corollary}[Restriction does not preserve action stability]\label{cor:failure}
Stability of action cotensors need not survive restriction along a
formal functor, even when its base map is the identity and both formal
categories have stable finite composition paths. This failure occurs
for a distributor on one actual bicomodule, whose restricted action
boundary has no contextual atomic representative.
\end{corollary}
\begin{proof}
Use the functor $\mathsf F:\mathsf A\to\mathsf B$ of
Corollary~\ref{cor:triangular-functor}. Lemma~\ref{lem:triangular} gives
stable composition paths for both categories. The $\mathsf B$-action
of Proposition~\ref{prop:stable-action} has a stable action cotensor;
its restriction is the $\mathsf A$-action of
Proposition~\ref{prop:restricted-action}. Lemma~\ref{lem:unstable-cotensor}
and Theorem~\ref{thm:strong} give the asserted instability and
nonrepresentability.
\end{proof}

\begin{remark}[A prime-indexed family]
Replacing $2,4$ by $p,p^2$ gives the same construction for every prime
$p$: use $\rho_H(b)=b\otimes g+pa\otimes x$, $N=\Z/p^2$, and
$R=\Z/p$. The problematic summand is $N[p]b$, its generator has
cone image $pb\otimes v$, and the reduced comparison again has matrix
$\operatorname{diag}(1,1,1,0)$. The proofs above apply with these
substitutions.
\end{remark}

\begin{thebibliography}{9}

\bibitem{Aguiar}
Marcelo Aguiar,
\emph{Internal categories and quantum groups},
Ph.D. thesis, Cornell University, 1997.
\url{https://pi.math.cornell.edu/~maguiar/thesis2.pdf}.

\bibitem{CS}
G.~S.~H. Cruttwell and Michael~A. Shulman,
\emph{A unified framework for generalized multicategories},
Theory and Applications of Categories \textbf{24} (2010), no.~21,
580--655.
\url{https://arxiv.org/abs/0907.2460}.

\end{thebibliography}
\end{document}
